\section{Limitations, non-goals, and open problems}
\label{sec:limitations}

The preceding sections describe what \Flux does, does not yet do, and is specified to do. This
section closes the description with three different kinds of boundary, and the three should not be
confused with one another. First, what the network's own roadmap states it will deliberately not
build, and why — a set of self-imposed limits on scope, not a gap in the argument. Second, what this
paper itself could not establish from the evidence available to it — measurements never taken,
provenance that cannot be pinned to a line of code, mechanisms documented nowhere in the assigned
repositories. Third, what remains genuinely unsolved: properties argued not to hold in earlier
sections, and research problems for which this paper states what a solution would have to look like
without supplying one. None of the three categories is filler; each is load-bearing for how a reader
should weigh the rest of the document.

\subsection{Stated non-goals}

The public roadmap carries its own section titled ``What we are not doing,'' positioned as a
top-level section immediately after the roadmap's discussion of supply and immediately before its
account of how it reports progress — not a subsection of anything else
\srcdoc{flux-roadmap/public/flux-roadmap-public.md:203}. It is reproduced below in full, item by
item, with its own stated reasoning and nothing added.

\begin{enumerate}
\item \textbf{``Reopening base emission upward.''} Reasoning given: ``\PNR v2 operates on revenue,
  not issuance.'' (\planned; \srcroadmap{What we are not doing, item 1; flux-roadmap-public.md:205})
\item \textbf{``Making privacy mandatory.''} Reasoning given: ``The transparent layer stays
  permanently.'' (\planned; \srcroadmap{What we are not doing, item 2; flux-roadmap-public.md:206})
\item \textbf{``Reselling operators' bandwidth or IP addresses'' for scraping.} Reasoning given in
  full: ``It is a real revenue idea and other networks do it. It exposes operators to risk they did
  not agree to. The answer is no, permanently.'' (\planned;
  \srcroadmap{What we are not doing, item 3; flux-roadmap-public.md:207})
\item \textbf{``\,`VPS on every node' in this cycle.''} Reasoning given: ``home connections cannot
  serve it reliably. The datacenter tier is what we can do properly now.'' (\planned;
  \srcroadmap{What we are not doing, item 4; flux-roadmap-public.md:208})
\item \textbf{``Adding new parallel-asset chains.''} Reasoning given: ``while we consolidate the ones
  we have.'' (\planned; \srcroadmap{What we are not doing, item 5; flux-roadmap-public.md:209})
\item \textbf{``Marketing a funnel we have not fixed.''} Reasoning given: ``The interface work comes
  first.'' (\planned; \srcroadmap{What we are not doing, item 6; flux-roadmap-public.md:210})
\end{enumerate}

No other list of non-goals appears anywhere else in the public roadmap corpus
\srcdoc{flux-roadmap/public/timeline.html; flux-roadmap/public/flux-roadmap-web.html}. A publicly
dated commitment not to do six specific things, each with its own stated reason, is unusual in the
systems-and-token-network literature surveyed for this paper, and among everything read for it, it is
the single piece of evidence that most directly argues for engineering judgment rather than messaging
setting the roadmap \srcinferred.

\subsection{Limitations of this paper}

Every entry below is recorded in \srcdoc{plan/gaps.md} as a claim the drafting process considered and
could not support, together with what the paper does instead of asserting it.

\paragraph{No realised measurement.} Two limitations are limits of public data, not of effort. First,
no public endpoint reports bandwidth, latency, or realised utilisation of the deployed fleet — only
requested resources from application specifications are visible. \S\ref{sec:evaluation} reports committed capacity
(8,972 vCPU, 17.1\,TiB RAM, 160\,TB disk) and states plainly that realised utilisation is not
measurable from public data, rather than estimating it (\srcdoc{plan/gaps.md, G13}). Second, operator
concentration is measurable only through the identities public chain data exposes. The owner-linked
proxy is \texttt{payment\_address}; \texttt{pubkey} is the node's VPS message-signing key and not the
owner key a vote would carry \srccode{fluxd/\allowbreak{}src/\allowbreak{}fluxnode/\allowbreak{}fluxnode.h}{154}. All such figures are lower
bounds, since every distinct identity is counted as a distinct party. §23 and \S\ref{sec:evaluation} state the figures as lower bounds
explicitly, and this section restates the true concentration measurement, by party rather than by
key, as an open problem below (\srcdoc{plan/gaps.md, G12}).

\paragraph{No competitive price comparison.} No pricing data for hyperscalers or peer networks was
gathered, and constructing a comparison would itself be exactly the kind of price claim the paper's
register constraint forbids. \S\ref{sec:related-work} compares architectures against related systems, never prices
(\srcdoc{plan/gaps.md, G14}).

\paragraph{Provenance the corpus could not settle.} Four specific facts resisted a citable,
unambiguous source when drafting began; two have since closed. The upstream \fluxd fork point cannot
be read off a version string, but is now bracketed from consensus evidence --- at or after
\texttt{zcashd} v2.1.0 (Blossom) and before v2.1.2 (Heartwood) --- rather than inferred from a
release-notes directory listing (§20; \srcdoc{plan/gaps.md, G15, closed}). The
\PoN activation height $\hpon = 2{,}020{,}000$ is a consensus fact, but the commonly cited activation
date does not appear anywhere in source; the nearest in-repository evidence is a checkpoint timestamp
that brackets, without establishing, the date, so the height is cited from code and the date only as
measured from a block explorer with a retrieval date (\srcdoc{plan/gaps.md, G16}). The Kadena parallel
asset's claim and payout mechanism was absent from the repository named for it, which contains no
code, and was located on a second pass in \texttt{fusion} and \texttt{flux-kda}; §6 now cites the
claim path and the single-key treasury transfer, and what remains undocumented is how that treasury
is replenished (\srcdoc{plan/gaps.md, G17, closed}).
And whether \FluxOS actually consumes the \texttt{flux-spec} schema could not be determined from
within \texttt{flux-spec} itself, which publishes no package and exports nothing, so §10 describes the
two artifacts as separate rather than presenting them as one specification
(\srcdoc{plan/gaps.md, G18}).

\paragraph{Claims tested and rejected.} Eleven further claims were considered during drafting and
found unsupportable; each is recorded in \srcdoc{plan/gaps.md} with the claim, why it fails, and what
the relevant section states instead. Four of them — that nodes vouch for what they serve, that
benchmarking establishes hardware fitness, that \PoN inherits Nakamoto-style reorg cost, and that
placement is deterministic — restate system properties argued fully in earlier sections and are not
re-argued here; they appear again, as properties rather than as refused claims, in the list of system
properties below (\srcdoc{plan/gaps.md, G1--G4}). The remaining seven are stated once, here, and nowhere else in
this paper: that the deployment assistant never signs is false on the default server-side path, so §19
states the property per signing path rather than as one property of the system
(\srcdoc{plan/gaps.md, G5}); that supply reaches a terminal value is false, so §5 gives the supply
function as unbounded with a vanishing rate and states its tail explicitly
(\srcdoc{plan/gaps.md, G6}); that \Flux implements a proof-of-useful-work mechanism is false — the
name survives only as a module name with no work-submission, verification, or reward protocol — so §15
describes the compute-rental marketplace as what it is and never invokes the term
(\srcdoc{plan/gaps.md, G7}); that GPUs are hot-attachable is unsupported by any code in the four
GPU-adjacent repositories, so §15 states whole-device passthrough as the deployed mechanism and tags
hot-attach \vision{} (\srcdoc{plan/gaps.md, G8}); that shielded supply is fully publicly verifiable is
weakened by turnstile enforcement being compiled in but disabled, so §20 states pool accounting as it
is and names turnstile activation as the condition that would make the claim true
(\srcdoc{plan/gaps.md, G9}); that application data survives a payment lapse is unsupported because no
grace, suspension, or eviction state machine exists in code, so §16 specifies one as \planned{} with
open parameters and states the deployed baseline — expiry arithmetic with no grace — separately
(\srcdoc{plan/gaps.md, G10}); and that the network compensates a tenant for node failure is
unsupported because no refund or SLA-credit mechanism exists anywhere, so §16 states that the remedy on
offer is replacement, not refund (\srcdoc{plan/gaps.md, G11}).

\subsection{Properties that do not hold}

The following are limitations of the deployed and specified system, not of this paper's evidence.
Each has already been argued in the section cited; they are collected here rather than re-argued.

\begin{itemize}
\item Benchmark results are not cryptographically bound to the operator or the host that produced
  them (§4). \shipped; \srccode{fluxbench/\allowbreak{}src/\allowbreak{}chainparams.cpp}{46--71}.
\item Attestations are not bound to the machine that produced them, and no tenant-facing verification
  path exists (§14). \shipped; \srccode{fluxsas/\allowbreak{}src/\allowbreak{}api\_server.h}{411--481}.
\item \PoN fork choice is a block-count race, not cumulative work: every block carries a fixed
  $2^{40}$ chainwork contribution regardless of its difficulty target (§4). \shipped;
  \srccode{fluxd/\allowbreak{}src/\allowbreak{}pow.cpp}{284--290}.
\item A 2-of-4 emergency key set can produce blocks with no time or stall-detection gating (§4).
  The set is intended to move to 3-of-4 with fresh keys \planned, which changes the threshold and
  not the property: the gating is absent at either size.
  \shipped; \srccode{fluxd/\allowbreak{}src/\allowbreak{}emergencyblock.cpp}{183--191}.
\item Placement is opportunistic self-selection with a 90-second broadcast collision wait, not
  deterministic scheduling (§9). \shipped; \srccode{flux/\allowbreak{}ZelBack/\allowbreak{}src/\allowbreak{}services/\allowbreak{}appLifecycle/\allowbreak{}appSpawner.js}{54--65,77--787}.
\item \texttt{fluxos-\allowbreak{}network-\allowbreak{}policy} is unsigned and takes fleet-wide effect within 7--13 hours,
  changeable without a required review step (§12). \shipped; \srccode{fluxos-network-policy/\allowbreak{}README.md}{30-32};
  \srccode{fluxos-network-policy/\allowbreak{}.github/\allowbreak{}CODEOWNERS}{1-15}.
\item \texttt{api.\allowbreak{}runonflux.\allowbreak{}io} is the sole application-discovery source, with no on-chain or
  peer-to-peer fallback (§12). \shipped; \srccode{flux-domain-manager/\allowbreak{}src/\allowbreak{}services/\allowbreak{}flux/\allowbreak{}dataFetcher.js}{47-91}.
\item The shielded pool cannot be entered: transparent value has been blocked from joining it since
  height 835,554 (§20). \shipped; \srccode{fluxd/\allowbreak{}src/\allowbreak{}main.cpp}{1398--1408}.
\item As specified, \PNR would create no marginal incentive to host, because payment reaches every
  node through the \PoN rotation independent of what that node runs (§7). \planned;
  \srcintent{mech-pnr.md}.
\item Four owner identities (three under linkage) reach the moderation eviction threshold at the measured
  collateral distribution (§23). \shipped; \srcmeasured{api.runonflux.io/\allowbreak{}daemon/\allowbreak{}listzelnodes}{2026-09-03}.
\item A \Flux light client is not constructible on an external chain under the current header format,
  which forecloses light-client bridging (§22). \shipped, as a consequence of the header facts above;
  \srccode{fluxd/\allowbreak{}src/\allowbreak{}pow.cpp}{284--290}; \srcintent{mech-bridge.md}.
\item GPU resource requests are silently stripped, rather than queued or rejected, when a host is
  contended (§15). \shipped; established in §15.
\item \texttt{Flux-Shared-DB} acknowledges a write before it has replicated (§13). \shipped;
  \srccode{Flux-Shared-DB/\allowbreak{}ClusterOperator/\allowbreak{}Operator.js}{1265-1338};
  \srccode{Flux-Shared-DB/\allowbreak{}ClusterOperator/\allowbreak{}server.js}{1543-1544}.
\end{itemize}

\subsection{Open problems}

The following four items are classed \emph{Research} in the feasibility assessment: each names an
open problem with no known good answer under the constraints \Flux operates within, rather than a
build that merely has not happened yet \srcdoc{plan/feasibility.md}. Each is stated with what would
have to be true of a solution, not merely that a solution is hard to find.

\begin{openproblem}[Post-quantum migration of outputs the protocol cannot move]
A repository-wide search of \fluxd's source returns no post-quantum primitive of any kind
\srcdoc{plan/feasibility.md, row 8}. The roadmap states the realistic threat as store-now-decrypt-later
rather than live breakage, and frames migration as a multi-year effort, explicitly placing \Flux where
Bitcoin and Ethereum are today: ``The realistic threat is store-now-decrypt-later, not live breakage.
We publish a migration path on the NIST-standardised signature schemes. Migration itself is a
multi-year effort — this is where Bitcoin and Ethereum are too.''
\srcroadmap{flux-roadmap-public.md:148--149}. A solution would have to address two different
populations of value separately. Node collateral is a registered, enumerable set, so a migration
targeting it first is tractable in a way a general UTXO migration is not — that is a structural
advantage, not a solved problem, and the paper states it rather than claiming credit for it
\srcdoc{plan/feasibility.md, per-mechanism notes}. Value sitting in ordinary transparent or shielded
outputs the protocol has no authority to move is the harder case: a solution there would need either
a voluntary, incentivised migration window before a cryptographically relevant quantum adversary is
credible, or an accepted statement that some outputs will simply not migrate.
\end{openproblem}

\begin{openproblem}[Bounded agent spending on a UTXO chain without a new trusted party]
Script validation on \fluxd is a pure function of the spending transaction and its own inputs; it
cannot observe transaction history, so no script-level rule can enforce a rate limit or a cumulative
spending cap \srcdoc{plan/feasibility.md, row 13}. The deployable answer available today — fund a
dedicated key and let its balance be the budget — is honest but coarse: it bounds total exposure, not
rate, and it cannot be revoked partway through a spend. A solution that bounds spending more richly
than a funded sub-key would have to supply either a co-signer, which reintroduces a trusted party the
rest of this paper's delegation design (§17) is built to avoid, or new consensus state that lets a
script condition on cumulative spend over a window — a Flux L1 consensus change with its own audit and
activation cost. Progress here means stating explicitly which of those two costs a design accepts, not
finding a way to avoid both.
\end{openproblem}

\begin{openproblem}[Unlinkable payment for a specific deployment]
An application's on-chain address is itself a workload identifier, so a construction proving that a
payment funded \emph{that} address defeats its own purpose the moment it succeeds
\srcdoc{plan/feasibility.md, row 18}. Two further constraints, drawn from the formal §21 design,
narrow what a solution could look like. First, amount hiding is unsatisfiable for any payment typed
by \PNR, because the fee pool is a public consensus integer and the payment amount is recoverable
from it to within a few satoshi — a solution has to be scoped to exclude \PNR-typed payments from
the start, not attempt to hide their amount. Second, the anonymity set is small by construction: at a
measured renewal rate, a renewal payment is a singleton within its own block essentially always and
within its own day the majority of the time, so a solution built on payer/workload unlinkability alone
starts from a small crowd, not a large one. A solution would have to be either a zero-knowledge
statement and circuit proving membership in a set of valid, unconsumed entitlements without naming
which one, verified within the existing block-interval budget, or a blind-credential design with a
third-party renewer that moves the trust assumption rather than removing it — and it would have to say
which of those two it is, rather than presenting payment-key indirection alone as a solution, since
indirection does not close the linkage on its own.
\end{openproblem}

\begin{openproblem}[Censorship resistance under measured collateral concentration]
At the measured ownership distribution, four owner identities (distinct \texttt{payment\_address}
values) already clear the 25\% eviction threshold, falling to three once addresses linked by a shared
node key or a shared collateral funding transaction are merged; the largest single identity holds
10.71\% of collateral-weighted voting power across 237 Stratus nodes
\srcdoc{plan/decisions-needed.md, item 11}. No reweighting, threshold change, or remedy within the
moderation quorum's design space lowers that number: the best available reweighting moves the
threshold from three colluding identities to six while making the threshold itself up to 13.6$\times$
cheaper to buy, which is a worse trade, not a better one. A solution is therefore not a mechanism-design
exercise inside the quorum; it is one of two things. Either it is a construction, such as the two-chamber
rule this paper's design work recommends, that raises the cost of a \emph{surprise} capture — requiring
a coalition to be pre-positioned and visible on-chain in advance — without claiming to raise the cost of
capture itself, since no such claim would be true; or it is an independent remeasurement of concentration
keyed to collateral-spending scripts rather than to node-signing keys, which could show the true
party-level concentration is lower than this lower bound suggests, and which has not been done. Progress
on this problem means one of those two things happening, not a better-tuned threshold.
\end{openproblem}

\subsection{Open decisions}

\srcdoc{plan/decisions-needed.md} records roughly eighty parameters that were open when drafting
began across the mechanisms this paper specifies. Reproducing all of them here would not be a
limitations section; it would be an appendix. Eighteen rounds of the design intake
\srcintent{08-answers-round-2.md} have since settled most of them, each in the section that uses the
value; the table below gives what is still open --- decisions no round has ruled on, and values this
paper proposed that the author has not confirmed, which round 18 requires to be marked as
recommendations rather than settled --- ordered by the section that carries each. Where the source does not
name who is responsible for a decision, the table says so rather than inventing a role.

%% longtable rather than table[H]: the previous [H] tabular ran 164pt past the text block (log:
%% overfull vbox), so its last row was cut and its caption never rendered. Same header pattern as
%% Appendix B. Re-cut around what is still open: round 18, G19.
\begingroup
\small
\begin{longtable}{@{}p{0.42\linewidth}p{0.28\linewidth}p{0.18\linewidth}@{}}
\caption{The decisions still open after rounds 2--18 of the design intake, ordered by the section
that carries each: items no round has ruled on, and values this paper proposed that the author has
not confirmed, each marked in place as \emph{Proposed by this paper}. The settled items this table
formerly listed are recorded in the sections that apply them.}
\label{tab:open-decisions}\\
\toprule
\textbf{Decision} & \textbf{Who decides} & \textbf{Blocks} \\
\midrule
\endfirsthead
\multicolumn{3}{l}{\small\itshape (Table~\ref{tab:open-decisions} continued)}\\
\toprule
\textbf{Decision} & \textbf{Who decides} & \textbf{Blocks} \\
\midrule
\endhead
\bottomrule
\endfoot
Emergency-block governance: the threshold moves to 3-of-4 with fresh keys, but whether the missing time or stall-detection gating is itself remedied
(time-lock, stall detection, retirement after \PoN stabilisation) & the author, explicitly & §4, §8, \S\ref{sec:migration} \\
\addlinespace
\PNR settlement: the four parameters still open after the author's confirmation of
\Cref{tab:pnr-recommended} --- O4 (fee floor and spam resistance), O8 (which payment channels must
settle through $\mathsf{sink}$, the load-bearing one), O14 (the identity of $\mathsf{sink}$) and O19
(the tier partition basis, to confirm only) & not named in source & §7, §16 \\
\addlinespace
The transformation per source version for the 322 applications on \specv{2}--\specv{7} under forced
automatic migration at $\hstar = \num{3050000}$; none is stated in any artefact read for this paper &
not named in source & §10, \S\ref{sec:migration} \\
\addlinespace
\texttt{quorum\allowbreak{}Grant\allowbreak{}Activation\allowbreak{}Height}: the observed \specv{9}-stability condition that triggers it and
the resulting height, deliberately decoupled from $\hstar$ \srcintent{08-answers-round-2.md, round 17}
& not named in source & §10 \\
\addlinespace
Bonded attestation network parameters --- $n = 9$, $k = 6$, $\Delta_{\mathrm{chal}} = \num{2880}$
blocks, full slashing, $\beta_i$ scaled to the value secured, re-checkable claims only at launch ---
are this paper's recommendations, unconfirmed & the author & §8, §14 \\
\addlinespace
Delegation certificate values --- the allowed-key set as a seventh scope component, a funded
allowance with no co-signer, no destination constraint, a re-issuance window of \num{8640} blocks,
digest resolution in the tooling --- are this paper's recommendations, unconfirmed & the author &
§17, §18 \\
\addlinespace
Two signing-path architecture: describe the local-signing and operator-signing paths, as §19 does, or
scope §19 to local signing alone; listed as still to be discussed
\srcintent{08-answers-round-2.md, round 2} and not returned to since & the author (editorial framing
decision) & §17, §18, §19 \\
\addlinespace
ZIP-209 turnstile enforcement with \Flux-correct constants for the drain-only window: recommended by
this paper; round 11 confirmed the window's heights, not the switch
\srcintent{08-answers-round-2.md, round 18} & the author & §5, §20, §21 \\
\addlinespace
Bridge: the checkable renunciation criteria (twelve months, \num{10000000}\flux{} of cumulative
volume, zero defect-fixing upgrades), the delaying implementation of never-refused burns, and direct
calldata notice to legacy holders are this paper's proposals; settled beneath them are the
audit-plus-operating-record shape (round 12) and the \num{10000000}\flux{} per-chain daily ceiling
(round 7) & the author & §22 \\
\addlinespace
Contract-held legacy balances on Ethereum and BSC: the set of LP, lending and bridge positions and
their per-protocol resolution to beneficial holders, to be published before $\hsnap$; the treatment
is settled \srcintent{08-answers-round-2.md, round 9} and the enumeration is not measured in this
corpus & Foundation and exchange partners & §22 \\
\addlinespace
Moderation category taxonomy: the enumeration of reportable categories, whose scope is ruled a
network-level content veto \srcintent{08-answers-round-2.md, round 9} and whose members are undecided
& not named in source & §23 \\
\end{longtable}
\endgroup

\subsection{What the next revision must do}

Some of this has closed on a schedule already stated elsewhere in this paper: the shielded pools'
sunset height and the $\specv{9}$ enforcement height are both settled (§20, §10), the fork point is
bracketed from consensus evidence (§20), and the activation-date provenance gap is a bounded task
— a corroborated explorer citation — rather than an open problem. Some of it
is not expected to close by the next revision, and should not be presented as though it will: the four
Research items above are open because the field does not have an answer, not because this paper did
not look hard enough, and a future draft that quietly drops one of them without a cited construction
should be read as having weakened its claim, not strengthened it. The standard by which a reader should
judge the next version of this document is not whether its count of open parameters and
\emph{not established} notes reaches zero. It is whether every claim that moved from \planned{} to \shipped{} did so because a
\texttt{[code:]} citation now exists for it, whether every \srcinferred{} claim that remains is still
visibly hedged, and whether the non-goals list above is still true.

%% Drain this section's float queue before the next begins.
\FloatBarrier
